\section*{Chapter \ref{ch:b3:current-account} --- The Current Account}

\begin{solution}{exo:b3:current-account:1}
In 2006, at 1.5~\% of world GDP. That year China's surplus was 0.45~\% of world GDP,
Japan's 0.34~\% and Germany's 0.32~\%; all other countries together had 0.58~\%, and the
sum over all countries was $+0.14~\%$, not zero. % figdata: ca_world_points
\end{solution}

\begin{solution}{exo:b3:current-account:2}
(a) \omterm{def:b3:current-account:accounts}{Current account}: export of goods $+100$; \omterm{def:b3:current-account:accounts}{financial account}: acquisition of a
foreign asset (the deposit) $+100$. (b) Primary income $+10$; a foreign asset (the
dividends, held abroad) $+10$. (c) \omterm{def:b3:current-account:accounts}{Financial account} only: incurrence of a liability
(the bonds) $+50$, acquisition of an asset (the kests received) $+50$; net zero.
(d) Secondary income $-5$; foreign assets $-5$. The \omterm{def:b3:current-account:accounts}{current account} is $100 + 10 - 5 =
105$, matched by net acquisition of foreign assets of $100 + 10 + 50 - 50 - 5 = 105$.
\end{solution}

\begin{solution}{exo:b3:current-account:3}
$CA = S - I = 22 - 25 = -3~\%$ of GDP. National saving falls by 2 points, to 20~\%, so
with investment unchanged the deficit widens to 5~\% of GDP.
\end{solution}

\begin{solution}{exo:b3:current-account:4}
Smoothing predicts that the \omterm{def:b3:current-account:accounts}{current account} falls when output is temporarily low
and rises when it is temporarily high. A series in which consumption falls one
for one with a temporary fall in output, and the \omterm{def:b3:current-account:accounts}{current account} does not move
--- or moves toward surplus in the slump --- contradicts it.
\end{solution}

\begin{solution}{exo:b3:current-account:5}
$\beta(1 + r) = 1$, so $C_1 = C_2 = (90 + 100/1.05)/(1 + 1/1.05) = 185.238/1.952 = 94.88$. $CA_1 = 90 -
94.88 = -4.88$; in period 2 net exports are $100 - 94.88 = 5.12$, interest paid $0.05 \times 4.88
= 0.24$, so $CA_2 = +4.88$: the debt is repaid. With $Y_1 = Y_2 = 90$, $C_1 = C_2 = 90$ and the
\omterm{def:b3:current-account:accounts}{current account} is zero in both periods.
\end{solution}

\begin{solution}{exo:b3:current-account:6}
$nx = 0.6 \times (0.04 - 0.03)/1.03 = 0.58~\%$ of GDP. With growth of 1~\%: $0.6 \times 0.03/1.01 =
1.78~\%$, three times as much.
\end{solution}

\begin{solution}{exo:b3:current-account:7}
The tested solution is \code{ex07_feldstein_horioka.py}, in the chapter's
\code{exercises/} folder. Slopes 0.129 (s.e.\ 0.063, 82 countries) for 1990--1999, 0.002
(0.048, 107) for 2000--2011 and 0.437 (0.053, 124) for 2012--2024. % output: ex07_feldstein_horioka
Not stable: the coefficient moves from about zero to 0.44 between periods, far
more than its standard errors. The samples differ (more countries report later),
and saving and investment both respond to common shocks such as commodity booms;
the slope measures capital mobility only under strong assumptions.
\end{solution}

\begin{solution}{exo:b3:current-account:8}
A permanent fall lowers permanent income by the same proportion as current
income: consumption adjusts at once and the \omterm{def:b3:current-account:accounts}{current account} barely moves. If the
fall is expected to reverse in two years, it is temporary: the country borrows to
keep consumption up, and the \omterm{def:b3:current-account:accounts}{current account} moves into deficit during the two
years.
\end{solution}

\begin{solution}{exo:b3:current-account:9}
(1) Expected returns on comparable assets equalised, for instance covered
interest parity holding with tiny deviations; (2) domestic investment rates
unrelated to domestic saving rates across countries. Others: consumption growth
highly correlated across countries (risk sharing); large gross flows in response to
small return differences.
\end{solution}

\begin{solution}{exo:b3:current-account:10}
$7.18^{1.5} = 19.2$ with $\alpha = 0.4$; $7.18^2 = 51.6$ with $\alpha = 1/3$. With human capital,
$\mathrm{MPK} = \alpha A^{1/\alpha}y^{(\alpha - 1)/\alpha}h^{(1 - \alpha)/\alpha}$, so the ratio is $(7.18/2.5)^{1.5} = 2.872^{1.5}
= 4.87$. Differences in human capital remove three quarters of the predicted gap.
\end{solution}

\begin{solution}{exo:b3:current-account:11}
$6/0.25 = 24~\%$ real depreciation. Through spending alone, imports must fall by 6~\% of
GDP: $0.3\,\Delta E = 6$, so spending falls by 20~\% of GDP --- why \omterm{def:b3:current-account:sudden-stop}{sudden stops} come
with deep recessions, and why the real exchange rate does most of the work when it
can.
\end{solution}

\begin{solution}{exo:b3:current-account:12}
The tested solution is \code{ex12_reversals.py}, in the chapter's \code{exercises/}
folder. 273 reversals in 118 countries; the \omterm{def:b3:current-account:accounts}{current account} averages $-12.3$, $-8.7$
and $-3.2~\%$ of GDP in years $t - 1$, $t$, $t + 1$; growth averages 4.2, 3.3, 3.6 and
3.0~\% in years $t - 2$ to $t + 1$. % output: ex12_reversals
On average growth slows by about a point before the reversal and does not
collapse: the average mixes forced reversals with deficits that closed because aid
or an investment project ended. The rule picks out large deficits (12~\% of GDP on
average) in small economies; a \omterm{def:b3:current-account:sudden-stop}{sudden stop} needs data on the inflows themselves.
\end{solution}

\begin{solution}{pb:b3:current-account:borrowing}
\textbf{1.} $CA = NX - \text{interest} = -278 - 1{,}786.5 = -2{,}064.5$ million dalens. % figdata: borrow_points, borrow_stats

\textbf{2.} 2025: net exports $+1{,}113$, but interest of 7{,}855 makes the \omterm{def:b3:current-account:accounts}{current
account} $-6{,}742$. (In 2008 and 2010 the trade surplus exceeded the interest bill.)

\textbf{3.} $-6{,}742/393{,}264 = -1.71~\%$ of GDP, financed by new net borrowing from
abroad: the debt rose by 6{,}742.

\textbf{4.} $63{,}053.3 - 64{,}908.7 = -1{,}855.4 = -CA_{2010}$.

\textbf{5.} Current-account deficits (trade deficits and the interest bill) raise
the debt; growth of GDP below the interest rate raises the ratio for a given debt.
(Valuation changes are a third channel, excluded here.) In 2020 both worked: a
deficit of 2.8~\% of GDP and nominal GDP down 3.2~\% (279{,}869 against 289{,}032).

\textbf{6.} Half of 146{,}511 is 73{,}256 in kests; its dalen value rises by 7{,}326, or
1.86~\% of 2025 GDP: the NIIP falls by that amount.

\textbf{7.} $C = (96.7 + 100/1.03)/(1 + 1/1.03) = 193.787/1.971 = 98.33$ in both periods.

\textbf{8.} $96.7 - 98.33 = -1.63~\%$ of normal output ($3.3/2.03$).

\textbf{9.} Brenland's \omterm{def:b3:current-account:accounts}{current account} went from $+0.50~\%$ of GDP in 2008 to $-1.71~\%$ in
2009, a fall of 2.21 points, somewhat more than the model's 1.63.

\textbf{10.} No change in the \omterm{def:b3:current-account:accounts}{current account}: consumption falls by the full 3.3.

\textbf{11.} The model has no investment, and investment usually falls in a
recession, which reduces the deficit; it also takes the shock to be known as
temporary and borrowing to be free at a fixed rate.

\textbf{12.} Real GDP fell by 4.85~\%; a temporary shock predicts a deficit widening
by about half of it, 2.4 points. The deficit widened from 1.67~\% to 2.77~\% of GDP, by
1.1 points: in the right direction, but less than full smoothing.

\textbf{13.} $d_t = d_{t-1}(1 + i_t)/(1 + g_t) - nx_t$.

\textbf{14.} $nx = 0.3726 \times (0.0590 - 0.0476)/1.0476 = 0.40~\%$ of GDP.

\textbf{15.} $0.28 < 0.40$: the ratio rises, by roughly 0.12 points of GDP a year at
these rates.

\textbf{16.} $0.3726 \times (0.04 - 0.0476)/1.0476 = -0.27~\%$: with $r < g$, a trade deficit of
0.27~\% of GDP keeps the ratio constant.

\textbf{17.} No: it requires only that the debt not grow faster than the interest
rate for ever. The debt can be rolled over indefinitely, provided the present value
of future trade balances covers it.

\textbf{18.} Repayment depends on the willingness to pay as well as the ability,
and there is no court to enforce it on a country; if each lender fears the others
will stop, a run can be self-fulfilling, since higher rates make the debt harder to
carry.

\textbf{19.} A \omterm{def:b3:current-account:accounts}{current account} of at least zero: no new net borrowing.

\textbf{20.} Net exports must cover the interest bill, $7{,}855/393{,}264 = 2.0~\%$ of GDP; the
swing from 0.28~\% is 1.71 points, the 2025 current-account deficit.

\textbf{21.} $1.71/0.3 = 5.7~\%$.

\textbf{22.} Brenland's external debt stabilises at 37~\% of GDP with a trade surplus
of 0.40~\% of GDP; a \omterm{def:b3:current-account:sudden-stop}{sudden stop} would require 2.0~\%.
\end{solution}
